%% TeX Live 安装说明（墨灵TeX 附带）— 详细版
\documentclass[UTF8,11pt]{article}
\usepackage[margin=1.8cm,top=1.6cm]{geometry}
\usepackage{ctex}
\usepackage{enumitem}
\usepackage{hyperref}
\usepackage{xcolor}
\usepackage{graphicx}
\usepackage{booktabs}

\setCJKmainfont{SimSun}
\setCJKsansfont{Microsoft YaHei}
\setCJKmonofont{FangSong}

\definecolor{accent}{HTML}{2563EB}
\definecolor{warn}{HTML}{B45309}
\definecolor{ok}{HTML}{047857}
\pagestyle{plain}
\setlength{\parindent}{0pt}
\setlist{itemsep=0.15em,topsep=0.25em,parsep=0pt}
\hypersetup{colorlinks=true,linkcolor=accent,urlcolor=accent}

\title{\textcolor{accent}{\sffamily\bfseries TeX Live 安装与配置指南}\\[0.4em]
\large 墨灵TeX 附带说明 · Windows}
\author{}
\date{}

\begin{document}
\maketitle
\vspace{-1.6em}

\begin{center}
\colorbox{accent!8}{\parbox{0.96\textwidth}{%
\textbf{适用环境：}Windows 10 / 11（64 位）\\
\textbf{推荐版本：}TeX Live 2024（2022 / 2023 亦可）\\
\textbf{安装完成后：}墨灵TeX 会自动检测 \texttt{bin} 目录；未检测到时可在「设置 $\rightarrow$ 编译」中手动填写。
}}
\end{center}

\vspace{0.6em}
\section*{一、为什么需要 TeX Live}
墨灵TeX 的\textbf{编辑}与\textbf{PDF 预览}可以独立使用，但要把 \texttt{.tex} 编译成 PDF，必须依赖本机 TeX 发行版。

\begin{itemize}
  \item 编译引擎：默认 \textbf{XeLaTeX}（中文友好），期刊/部分模板可能用 \textbf{pdfLaTeX}
  \item 宏包管理：ctex、amsmath、geometry、hyperref 等均由发行版提供
  \item 首次安装约需磁盘 \textbf{5--8 GB}（完整方案），网络较差时请预留更久时间
\end{itemize}

\section*{二、下载安装包}

\subsection*{方式 A：官方安装器（推荐新手）}
\begin{enumerate}
  \item 打开官网：\url{https://tug.org/texlive/}
  \item 进入 \textbf{Getting TeX Live} 页面
  \item Windows 用户下载 \textbf{install-tl-windows.exe}（在线安装器，体积较小）
  \item 若网络较慢，可改用 \textbf{install-tl-w64.exe} 或从页面中的 CTAN 镜像站获取 ISO
\end{enumerate}

\subsection*{方式 B：完整 ISO（适合离线/校园网）}
\begin{enumerate}
  \item 在 \textbf{Getting TeX Live} 页选择 \textbf{on DVD / ISO image}
  \item 从国内镜像下载完整 ISO（清华、中科大等），体积约数 GB
  \item 用资源管理器挂载 ISO（Windows 10/11 可直接双击挂载）
\end{enumerate}

\textcolor{warn}{\textbf{提示：}请务必从 TUG 或可信镜像下载，避免第三方捆绑安装包。}

\section*{三、Windows 安装步骤（逐步）}

\begin{enumerate}
  \item 右键安装程序 $\rightarrow$ \textbf{以管理员身份运行}（可减少权限问题）
  \item 安装语言可选 \textbf{English} 或简体中文界面
  \item 安装类型：
    \begin{itemize}
      \item \textbf{Advanced / Custom}（推荐）：可核对路径与组件
      \item \textbf{Default}：适合不想改配置的用户
    \end{itemize}
  \item \textbf{安装目录}建议保持默认，例如：
    \begin{quote}
      \texttt{C:\textbackslash texlive\textbackslash 2024}
    \end{quote}
    自定义时请记住最终 \texttt{bin\textbackslash windows} 路径，稍后要在墨灵里使用。
  \item 组件选择：
    \begin{itemize}
      \item 保持默认全选，即可包含中文支持（XeTeX / LuaTeX / ctex 相关包）
      \item 磁盘紧张时可取消部分文档与多语言包，但务必保留 \textbf{XeTeX}、\textbf{LuaTeX}、\textbf{宏包集合}
    \end{itemize}
  \item 点击 \textbf{Install}，等待进度条走完（约 20--60 分钟，视网速而定）
  \item 结束页若提示「安装完成」即可关闭安装器
\end{enumerate}

\section*{四、验证安装是否成功}
按 \textbf{Win+R}，输入 \texttt{cmd} 回车，依次尝试：

\begin{enumerate}
  \item 输入 \texttt{xelatex -version}
  \item 若提示「不是内部或外部命令」，说明 PATH 未生效：
    \begin{itemize}
      \item 关闭并重新打开终端窗口再试
      \item 或重启电脑
    \end{itemize}
  \item 若显示版本号（如 TeX Live 2024），表示命令行已可用
\end{enumerate}

也可直接打开墨灵TeX，看状态栏/设置里是否自动识别到 TeX Live。

\section*{五、在墨灵TeX 中配置路径}

\begin{enumerate}
  \item 启动墨灵TeX，软件启动后约数秒会自动探测 TeX Live
  \item 若右上角/状态提示已就绪，可直接打开模板或文件按 \textbf{F5} 编译
  \item 若未检测到：
    \begin{enumerate}
      \item 菜单栏 \textbf{设置}（或命令面板搜「设置」）
      \item 切到 \textbf{编译} 分区
      \item 在 \textbf{TeX Live 路径} 中填入 bin 目录，例如：
\begin{quote}
\texttt{C:\textbackslash texlive\textbackslash 2024\textbackslash bin\textbackslash windows}
\end{quote}
      \item 保存设置，重新编译
    \end{enumerate}
\end{enumerate}

\section*{六、常用编译引擎选择}

\begin{center}
\begin{tabular}{@{}p{3.2cm}p{5.2cm}p{4.8cm}@{}}
\toprule
文档类型 & 推荐引擎 & 说明 \\
\midrule
中文论文 / 简历 / 幻灯片 & XeLaTeX & 支持系统中文字体（如宋体、雅黑） \\
国赛 CUMCM & XeLaTeX & 模板要求 xelatex \\
IEEE / Elsevier & pdfLaTeX & 英文期刊常见 \\
纯英文 article & 两者皆可 & 默认跟墨灵设置 \\
\bottomrule
\end{tabular}
\end{center}

\textcolor{ok}{\textbf{提示：}墨灵在编译时会根据 \texttt{documentclass} 自动纠正引擎；也可在设置中手动指定。}

\section*{七、常见问题排查}

\begin{enumerate}
  \item \textbf{找不到 xelatex / pdflatex}
  \begin{itemize}
    \item 确认路径是否指向 \texttt{...bin\textbackslash windows}（不是 texlive 根目录）
    \item 该目录下应能看到 \texttt{xelatex.exe}
  \end{itemize}
  \item \textbf{中文乱码或无法编译}
  \begin{itemize}
    \item 引擎务必用 XeLaTeX
    \item 文档类使用 ctexart/ctexbeamer 时不要用 pdfLaTeX
  \end{itemize}
  \item \textbf{缺少宏包（如 fontawesome5）}
  \begin{itemize}
    \item 完整 TeX Live 一般已含
    \item MiKTeX 会提示按需安装，允许即可
  \end{itemize}
  \item \textbf{杀毒/防火墙拦截}
  \begin{itemize}
    \item 首次编译可能被拦截，请将墨灵TeX 与 TeX Live 目录加入信任
  \end{itemize}
  \item \textbf{安装很慢或中断}
  \begin{itemize}
    \item 改用国内镜像，或改用 ISO 离线安装
    \item 中断后重新运行安装器，通常可续装
  \end{itemize}
\end{enumerate}

\section*{八、可选：MiKTeX 替代方案}
若不想安装体积很大的 TeX Live，可安装 \textbf{MiKTeX}（\url{https://miktex.org/}）：
\begin{itemize}
  \item 体积更小，缺包时可自动下载
  \item 安装后同样需要在墨灵「设置 $\rightarrow$ 编译」中指定其 \texttt{miktex\textbackslash bin} 路径
  \item 部分中文模板仍更推荐 TeX Live + XeLaTeX
\end{itemize}

\vspace{1em}
\begin{center}
\colorbox{accent!8}{\parbox{0.96\textwidth}{%
完成安装并配置路径后，可从起始页「浏览模板库」选择模板，或打开已有 \texttt{.tex} 按 \textbf{F5} 编译。
若仍失败，请查看底部日志面板中的完整错误信息。
}}
\end{center}

\end{document}
